% ============================================================
%  第二部分之三：非参数、机器学习与工程规则法
% ============================================================

\section{非参数、机器学习与工程规则法}

前面两章处理的是“有显式概率模型”的路线。本章覆盖另一半地图：不假设发射分布的非参数方法（聚类）、把识别外包给学习器的机器学习路线（监督与深度），以及工程上久经考验的规则法。三者不是替代关系，而是成本、可解释性与表达力上的不同取舍。

\subsection{聚类族：从 K-means 到 GMM}

聚类把每个时期的特征向量（收益率、波动率、相关性、流动性……）映射为一个状态标签，是最朴素的状态识别。

\begin{itemize}[itemsep=2pt]
\item \textbf{K-means}：硬聚类，最小化簇内平方距离。优点：快、无需分布假设；缺点：簇形状受欧氏距离支配、硬边界无法表达过渡期。一个重要联系：统计跳变模型在跳变惩罚 $\lambda = 0$ 时\emph{精确退化为} K-means~\cite{shu2024a}——这说明两者不是两个阵营，而是一个谱系的两端（第~\ref{sec:frontier}~节）。
\item \textbf{高斯混合（GMM）}：软聚类，输出后验概率。GMM 是“去掉时序结构”的 HMM——两者共享同样的发射分布假设，差别只在于是否使用转移矩阵。若在 GMM 上恢复 Markov 先验，就得到 HMM；若在 HMM 上抹掉转移结构（$a_{ij}$ 与 $i$ 无关），就退化回 GMM。这个互推关系是理解整张方法地图的枢纽。
\item \textbf{谱聚类 / DBSCAN}：前者适合非凸簇结构与图结构（如相关性网络聚类）；后者把低密度点视为噪声，适合把“危机时期”当作异常簇的建模视角。
\end{itemize}

\begin{remark}[聚类的真正难点是特征而不是算法]
实践中决定成败的从来不是“用 K-means 还是 GMM”，而是\emph{拿什么去聚类}：只有收益率的聚类在噪声里挣扎；加入已实现波动率、成交量、相关性均值、跳跃统计量（已实现方差中的跳变分解）后，簇结构才稳定。特征设计决定了状态的经济含金量，这一点在第~\ref{sec:engineering}~节剖析的实战体系里同样成立（其 L3 层用整整一组技术因子作为判别输入）。
\end{remark}

\subsection{监督路线：把识别转化为预测}

无监督方法只能回答“当下像什么”，无法直接回答“下一期会是什么”。监督路线把问题改写为预测问题：以无监督标签为 $y$、以可观测特征为 $x$，训练分类器。

Shu et al.~\cite{shu2024b} 的框架是这一路线的代表作：先用统计跳变模型生成历史状态标签（牛/熊），再以“资产自身特征 $+$ 跨资产宏观特征”训练梯度提升树分类器预测状态；预测结果进入 Markowitz~\cite{markowitz1952} 均值—方差优化。该文在 1991--2023 年十二类跨资产组合（全球股票、债券、房地产、商品指数）上报告了对最小方差、均值方差与朴素分散组合的一致改善。

\begin{remark}[标签质量决定天花板，而标签泄漏是最隐蔽的坑]
监督路线的上限由标签质量决定。若标签本身用\emph{全样本}平滑或离线分割生成（例如用全样本 Bai--Perron 或全样本 HMM 平滑概率），监督模型就在学习“未来信息塑造的标签”，任何回测绩效都不可信。正确做法与第~\ref{sec:validation}~节的检查清单一致：标签必须逐期用扩展窗口重新生成（walk-forward labeling）。
\end{remark}

\subsection{深度学习路线：能力与边界}

\subsubsection{自编码器：以重构误差为信号}

用正常时期数据训练自编码器，把“重构误差异常”当作结构变化或危机信号。优点是无需标签、可处理高维；缺点是阈值的先验性、以及对“缓慢漂移”的迟钝（重构误差对渐变不敏感，对突跳敏感——与变点检测的谱系一致）。

\subsubsection{可微状态的尝试：KASPER}

Oad et al.~\cite{oad2025} 的 KASPER 把“状态”做成可微结构：用 Gumbel-Softmax 实现可微的 Regime 分类层，配合对比损失（拉大类间表示距离）与正交正则，再用基于蒙特卡洛 Shapley 值的符号规则提取，输出每状态可读的决策规则。它的价值主张不在精度（其报告的 $R^2 = 0.89$、Sharpe $12.02$ 等指标应在“样本与设定依赖”的严格限定下解读），而在于展示了一条“深度模型 $+$ 可解释性”的工程路径。

\subsubsection{非参数谱方法：Hilbert--Huang 变换}

Luwang et al.~\cite{luwang2026} 用经验模态分解（EMD）驱动的 Hilbert--Huang 变换，从瞬时能量谱划分 Normal / High / Extreme 三类状态，再以变长 Markov 链估计各状态内的收益转移。其结论对工程有两点直接价值：
\begin{itemize}[itemsep=2pt]
\item \emph{状态差异体现在方差机制而非均值}：极端状态下尾部状态概率占优，而熵在“高波动”状态达到峰值——不确定性的峰值出现在中等波动区，而不是极端区；
\item \emph{成熟度影响状态恢复力}：发达市场在压力消退后更快回归常态，发展中市场在“正常”状态仍保留尾部依赖——这与“同一状态模型跨市场迁移需重新校准”的工程经验一致。
\end{itemize}

\subsection{强化学习：把状态接入决策}

Verma et al.~\cite{verma2026} 在 SPY/TLT/GLD 三资产上把三状态高斯 HMM（BIC 选择状态数）与表格 MDP 强化学习结合：状态判定交给 HMM，配置决策交给 RL 策略。2004--2025 数据的样本外结果（30\% 测试窗、1 日执行延迟）显示：HMM 条件化配置已超过被动基准，RL 版本在风险调整后表现最优、回撤显著更低，且行为可解释（每个状态对应明确的动作映射）。

\begin{remark}[RL 路线的诚实边界]
RL 在 Regime 配置中的合理性来自“策略 $=$ 状态相关的动作规则”这一结构先验；其限制同样清楚：奖励信号（收益）信噪比极低、样本量有限（金融时间序列只有一个现实路径）、状态判定误差会污染策略梯度。因此当前可行的方法是“小状态空间 $+$ 有限动作 $+$ 强先验”的保守版本，而非端到端大模型。
\end{remark}

\subsection{工程规则法：ADX、Hurst、谱分析与波动率分位}

规则法用少量指标与阈值直接划分趋势/震荡、高波/低波。它们不优雅，但透明、无需估计、没有收敛问题，且在工业界广泛使用。

\begin{itemize}[itemsep=2pt]
\item \textbf{ADX（平均趋向指数）}：用方向运动（$+DI$、$-DI$）与真实波幅构造趋势强度，$\mathrm{ADX} > 25$（或归一化后的等价阈值）被普遍视为“趋势状态”的粗糙判据。第~\ref{sec:engineering}~节的实战体系在 L3 层使用归一化 ADX（阈值 $0.25$）作为趋势确认信号。
\item \textbf{Hurst 指数}：度量长记忆性，$H > 0.5$ 偏向趋势持续、$H < 0.5$ 偏向均值回归。它更接近“策略适配性”而非“市场状态”的语言——本工作区在趋势适配性评估中即用 Hurst $+$ 频域方法 $+$ HMM 的组合。
\item \textbf{FFT / 谱分析}：把“周期性 vs 趋势性”的能量分布量化，适合识别显著的日内/周度模式与低频能量。
\item \textbf{波动率分位}：把当前已实现波动率映射到历史分位数。它是所有规则法中最稳健的单变量状态代理——因为波动率本身具有强持久性（第一部分的三大经验事实之一），分位数的排序信息几乎不依赖状态假设。
\end{itemize}

\begin{remark}[规则法的定位：低成本的“判别式”方法]
在第一章的三种刻画分类中，规则法属于典型的\emph{判别式}：给定特征直接输出标签，不产生概率、不承诺似然。它的两个后果需要被工程化处理：（1）阈值的先验性——阈值不是从数据里学来的，因此不同参数体系下“同一状态”的名义定义不同；（2）状态无概率——无法做“状态置信度”与“状态概率”下游接口。正因如此，实战系统常把规则法\emph{与}统计模型组合（投票、加权、级联），用两种不同的偏置来互相约束——这正是第~\ref{sec:engineering}~节体系设计中“HMM $+$ 规则法投票”的动机。
\end{remark}

\subsection{方法选择矩阵}

\begin{table}[htbp]
\centering
\footnotesize
\begin{tabular}{>{\raggedright\arraybackslash}p{2.7cm}>{\raggedright\arraybackslash}p{2.3cm}>{\raggedright\arraybackslash}p{2.1cm}>{\raggedright\arraybackslash}p{2.0cm}>{\raggedright\arraybackslash}p{2.1cm}}
\toprule
\textbf{方法族} & \textbf{数据要求} & \textbf{输出} & \textbf{时序结构} & \textbf{典型场景} \\
\midrule
K-means / 谱聚类 & 低 & 硬标签 & 无 & 探索性结构分析 \\
GMM & 低 & 软概率 & 无 & 快速状态画像 \\
HMM / MS & 中 & 滤波后验 & 完整 & 状态路由、风险预算 \\
HMM $+$ 深度特征 & 高 & 后验 & 完整 & 高维微观结构 \\
监督两阶段 & 高 & 预测概率 & 依赖标签 & 状态预测、配置 \\
规则法 & 极低 & 标签 & 无 & 快速筛查、投票组件 \\
\bottomrule
\end{tabular}
\caption{状态识别方法的工程选择矩阵}
\label{tab:method-matrix}
\end{table}

\keynote{聚类、监督学习、深度模型与规则法不是四个阵营，而是一条谱系上的四个位置：谱系的一端是“无分布假设 $+$ 无时序结构”，另一端是“完整似然 $+$ 完整时序”；选择位置的标准不是先进性，而是你对数据量与信噪比的诚实判断。}